\section{Certain Identities Involving Fibonacci Numbers}

\begin{exercise}
    Using induction on the positive integer $n$, establish the formulas
    \begin{enumerate}
        \item $u_1 + 2u_2 + 3u_3 + \dots + nu_n = (n+1)u_{n+2} - u_{n+4} + 2$;
        \item $u_2 + 2u_4 + 3u_6 + \dots + nu_{2n} = nu_{2n + 1} - u_{2n}$.\\
    \end{enumerate}
\end{exercise}

\begin{solution}
    \begin{enumerate}
        \item The case $n = 1$ holds since both $u_1$ and $(1+1)u_3 - u_5 + 2$ are equal to 1. Next, suppose that
        $$u_1 + 2u_2 + 3u_3 + \dots + ku_k = (k+1)u_{k+2} - u_{k+4} + 2$$
        for some $k \geq 1$, then adding $(k+1)u_{k+1}$ pn both sides gives us
        $$u_1 + 2u_2 + 3u_3 + \dots + (k+1)u_{k+1} = (k+1)(u_{k+2} + u_{k+1}) - u_{k+4} + 2.$$
        Now, notice that $(k+1)(u_{k+2} + u_{k+1}) = (k+1)u_{k+3}$ which lets us rewrite $$(k+1)(u_{k+2} + u_{k+1}) - u_{k+4} + 2 = (k+2)u_{k+3} - (u_{k+4} + u_{k+4}) + 2.$$
        It follows that,
        $$u_1 + 2u_2 + 3u_3 + \dots + (k+1)u_{k+1} = ((k+1) + 1)u_{(k+1)+2} - u_{(k+1) + 4} + 2,$$
        which shows that the case $n = k+1$ holds as well. Therefore, it holds for all $n \geq 1$ by induction.
        \item The case $n = 1$ holds since both $u_2$ and $1\cdot u_3 - u_2$ are equal to 1. Next, suppose that
        $$u_2 + 2u_4 + 3u_6 + \dots + ku_{2k} = ku_{2k + 1} - u_{2k}$$
        for some $k \geq 1$, then adding $(k+1)u_{2k+2}$ on both sides gives us
        $$u_2 + 2u_4 + 3u_6 + \dots + (k+1)u_{2(k+1)} = (k+1)u_{2k+2}+ ku_{2k + 1} - u_{2k}.$$
        Now, notice that 
        $$(k+1)u_{2k+2}+ ku_{2k + 1} - u_{2k} = (k+1)(u_{2k+2} + u_{2k+1}) - (u_{2k} + u_{2k + 1}) = (k+1)u_{2k + 3} - u_{2k + 2}.$$
        It follows that,
        $$u_2 + 2u_4 + 3u_6 + \dots + (k+1)u_{2(k+1)} = (k+1)u_{2(k+1) + 1} - u_{2(k+1)},$$
        which shows that the case $n = k+1$ holds as well. Therefore, it holds for all $n \geq 1$ by induction. \\
    \end{enumerate}
\end{solution}

\begin{exercise}
    \begin{enumerate}
        \item Show that the sum of the first $n$ Fibonacci numbers with odd indices is given by the formula
        $$u_1 + u_3 + u_5 + \dots + u_{2n-1} = u_{2n}.$$
        \hint{Add the equalities $u_1 = u_2$, $u_3 = u_4 - u_2$, $u_5 = u_6 - u_4$, \dots}
        \item Show that the sum of the first $n$ Fibonacci numbers with even indices is given by the formula
        $$u_2 + u_4 + u_6 + \dots + u_{2n} = u_{2n+1} -1.$$
        \hint{Apply part (a) in conjunction with identity (2)}
        \item Derive the following expression for the alternating sum of the first $n$ Fibonacci numbers:
        $$u_1 - u_2 + u_3 - u_4 + \dots + (-1)^{n+1}u_{n} = 1 + (-1)^{n+1}u_{n-1}.$$
    \end{enumerate}
\end{exercise}

\begin{solution}
    \begin{enumerate}
        \item Since $u_1 = u_2$, $u_3 = u_4 - u_2$, $u_5 = u_6 - u_4$, ..., $u_{2n - 1} = u_{2n} - u_{2n - 2}$, then we obtain the following telescoping sum:
        $$u_1 + u_3 + u_5 + \dots + u_{2n-1} = u_2 + (u_4 - u_2) + (u_6 - u_4) + \dots + (u_{2n} - u_{2n - 2}) = u_{2n}.$$
        \item If we let $S_1 = \sum_{k=1}^{2n}u_k$ and $S_2 = \sum_{k=1}^{n}u_{2k-1}$, then on one hand we have that $\sum_{k=1}^{n}u_{2k} = S_1 - S_2$, and on the other hand we have that $S_1 = u_{2n+2} - 1$ (identity (2)) and $S_2 = u_{2n}$ (part (a)). Therefore, we obtain
        $$u_2 + u_4 + u_6 + \dots + u_{2n} = u_{2n+2} - 1 - u_{2n} = u_{2n + 1} - 1.$$
        \item First, let's prove the case where $n$ is even. In that case, we let $S = \sum_{k=1}^{n}(-1)^{k+1}u_k$, $S_1 = \sum_{k=1}^{n}u_k$ and $S_2 = \sum_{k=1}^{n/2}u_{2k-1}$. On one hand, we have that $S + S_1 = 2S_2$, on the other hand we have that $S_1 = u_{n+2} - 1$ by identity (2) and $S_2 = u_n$ by part (a). Thus, $S = 2u_n - u_{n+2} + 1 = 1 - u_{n+1}$. Therefore, we have that
        $$u_1 - u_2 + u_3 - u_4 + \dots + (-1)^{n+1}u_{n} = 1 + (-1)^{n+1}u_{n-1}$$
        holds in that case since the sum of the left hand side is $S$. 
        
        Next, suppose that $n$ is odd, define $S$ and $S_1$ as before and define $S_2 = \sum_{k=1}^{(n+1)/2}u_{2k - 1}$. Again, we have that $S + S_1 = 2S_1$, and that $S_1 = u_{n+2} - 1$ and $S_2 = u_{n+1}$. It follows that $S = 2u_{n+1} - u_{n+2} + 1 = 1 + u_{n-1}$. Therefore, we have that
        $$u_1 - u_2 + u_3 - u_4 + \dots + (-1)^{n+1}u_{n} = 1 + (-1)^{n+1}u_{n-1}$$
        holds in that case, and thus, for all $n \geq 1$. \\
    \end{enumerate}
\end{solution}

\begin{exercise}
    From formula (1), deduce that
    $$u_{2n - 1} = u_n^2 + u_{n-1}^2, \qquad u_{2n} = u_{n+1}^2 - u_{n-1}^2, \quad \text{ with } n\geq 2.$$
\end{exercise}

\begin{solution}
    In formula (1), if we replace $n$ with $n-1$ and $m$ with $n$, we obtain 
    $$u_{2n - 1} = u_{n-1}^2 + u_n^2.$$
    Similarly, if we now replace $n$ with $n - 1$ and $m$ with $n +1$, we obtain
    $u_{2n} = u_nu_{n-1} + u_{n + 1}u_n$. Factorizing by $u_n$ and rewriting $u_n$ as $u_{n+1} - u_{n-1}$ gives us
    $$u_{2n} = (u_{n+1} - u_{n-1})(u_{n+1} + u_{n-1}) = u_{n+1}^2 - u_{n-1}^2.$$ \\
\end{solution}

\begin{exercise}
    Use the results of Problem 3 to obtain the following identities:
    \begin{enumerate}
        \item $u_{n+1}^2 + u_{n-2}^2 = 2u_{2n-1}; \hfill n \geq 3$
        \item $u_{n+2}^2 + u_{n-1}^2 = 2(u_n^2 + u_{n+1}^2). \hfill n \geq 2$\\
    \end{enumerate} 
\end{exercise}

\begin{solution}
    \begin{enumerate}
        \item By the Problem 3, we have that $u_{n+1}^2 - u_{n-1}^2 = u_{2n}$, $u_n^2 + u_{n-1}^2 = u_{2n-1}$ and $u_n^2 - u_{n-2}^2 = u_{2n-2}$. Adding these three equations together gives us
        $$u_{n+1}^2 + u_{n-2}^2 = (u_{n+1}^2 - u_{n-1}^2) + (u_n^2 + u_{n-1}^2) + (u_n^2 - u_{n-2}^2) = u_{2n} + u_{2n-1} + u_{2n-2} = 2u_{2n}.$$
        \item If we replace $n$ with $n + 1$ in part (a), we obtain
        $$u_{n+2}^2 + u_{n-1}^2 = 2u_{2(n+1)} = 2(u_n^2 + u_{n+1}^2).$$ \\
    \end{enumerate}
\end{solution}

\begin{exercise}
    Establish that the formula
    $$u_nu_{n-1} = u_n^2 - u_{n-1}^2 + (-1)^n$$
    holds for $n \geq 2$ and use this to conclude that the consecutive Fibonacci numbers are relatively prime. \\
\end{exercise}

\begin{solution}
    By formula (3), we have that
    $$u_n^2 - u_{n-1}^2 + (-1)^n = u_{n+1}u_{n-1} + (-1)^{n-1} - u_{n-1}^2 + (-1)^n.$$
    Now, notice that $(-1)^{n-1} + (-1)^n = 0$ and $u_{n+1}u_{n-1} - u_{n-1}^2 = u_{n-1}(u_{n+1} - u_{n-1}) = u_{n-1}u_n$. Therefore, we obtain the formula
    $$u_n^2 - u_{n-1}^2 + (-1)^n = u_{n-1}u_n$$
    which is the formula we wanted.
    
    Next, suppose that $u_{n-1}$ and $u_n$ have a nonzero divisor $d > 1$, then
    $$d^2\left(\left(\frac{u_n}{d}\right)^2 - \left(\frac{u_{n-1}}{d}\right)^2 - \frac{u_{n-1}u_n}{d^2}\right) = (-1)^{n+1}.$$
    But this implies that $d \mid (-1)^{n+1}$, and hence, that $d = \pm 1$. This is a contradiction with the fact that $d > 1$ so $\gcd(u_{n-1}, u_n) = 1$. \\
\end{solution}

\begin{exercise}
    Without resorting to induction, derive the following identities:
    \begin{enumerate}
        \item $u_{n+1}^2 - 4u_n u_{n-1} = u_{n-2}^2. \hfill n\geq 3$
        
        \hint{Start by squaring both $u_{n-2} = u_n - u_{n-1}$ and $u_{n+1} = u_n + u_{n-1}$}
        \item $u_{n+1}u_{n-1} - u_{n+2}u_{n-2} = 2(-1)^n. \hfill n\geq 3$
        
        \hint{Put $u_{n + 2} = u_{n+1} + u_n$, $u_{n-2} = u_n - u_{n-1}$, and use formula (3)}
        \item $u_n^2 - u_{n+2}u_{n-2} = (-1)^n. \hfill n\geq 3$.
        
        \hint{Mimic the proof of formula (3)}
        \item $u_n^2 - u_{n+3}u_{n-3} = 4(-1)^{n+1}. \hfill n\geq 4$
        \item $u_nu_{n+1}u_{n+3}u_{n+4} = u_{n+2}^4 - 1, \hfill n\geq 1$
        
        \hint{By part (c), $u_nu_{n+4} = u_{n+2}^2 + (-1)^{n+1}$, while by formula (3), $u_{n+1}u_{n+3} = u_{n+2}^2 + (-1)^{n+2}$}
        \item $u_{2n+2}u_{2n-1} - u_{2n}u_{2n+1} = 1. \hfill n \geq 1$ \\
    \end{enumerate}
\end{exercise}

\begin{solution}
    \begin{enumerate}
        \item Since $u_{n + 1}^2 = u_n^2 + 2u_n u_{n-1} + u_{n-1}^2$, then
        $$u_{n + 1}^2 - 4u_nu_{n-1} = u_n^2 - 2u_n u_{n-1} + u_{n-1}^2 = (u_n - u_{n-1})^2 = u_{n-2}^2.$$ 
        \item By formula 3, we have that $u_{n+1}u_{n-1} = u_n^2 + (-1)^n$. Hence,
        $$u_{n+1}u_{n-1} - u_{n+2}u_{n-2} = u_n^2 - u_{n+2}u_{n-2} + (-1)^n.$$
        Next, rewriting $u_{n+2}$ as $u_{n+1} + u_n$ gives us
        $$u_{n+1}u_{n-1} - u_{n+2}u_{n-2} = u_n^2 - u_{n+1}u_{n-2} - u_nu_{n-2} + (-1)^n.$$
        Again, formula (3) tells us that $u_n u_{n-2} = u_{n-1}^2 + (-1)^{n-1}$ so we get that 
        $$u_{n+1}u_{n-1} - u_{n+2}u_{n-2} = u_n^2 - u_{n+1}u_{n-2} - u_{n-1}^2 + 2(-1)^n.$$
        Now, we can notice that $u_n^2 - u_{n-1}^2 = (u_n - u_{n-1})(u_n + u_{n-1}) = u_{n-2}u_{n+1}$, and hence, $u_n^2 - u_{n+1}u_{n-2}- u_{n-1}^2 = 0$. Therefore, we obtain
        $$u_{n+1}u_{n-1} - u_{n+2}u_{n-2} = 2(-1)^n.$$
        \item By formula (3), we have $u_n^2 = u_{n+1}u_{n-1} - (-1)^n$. Thus, by part (b), we obtain
        $$u_n^2 - u_{n+2}u_{n-2} = u_{n+1}u_{n-1} - - u_{n+2}u_{n-2} - (-1)^n = (-1)^n.$$
        \item If we use the fact that $u_{n+3} = u_{n+2} + u_{n+1}$ and $u_{n-3} = u_{n-1} - u_{n-2}$, we get that
        $$u_n^2 - u_{n+3}u_{n-3} = u_n^2 - u_{n+2}u_{n-1} + u_{n+2}u_{n-2} - u_{n + 1}u_{n-1} + u_{n+1}u_{n-2}.$$
        By part (b), we get that 
        $$u_n^2 - u_{n+3}u_{n-3} = u_n^2 - u_{n+2}u_{n-1} + u_{n+1}u_{n-2} + 2(-1)^{n+1}.$$
        Next, replacing $u_{n+2}$ with $u_{n+1} + u_n$ gives us
        $$u_n^2 - u_{n+3}u_{n-3} = u_n^2 - u_{n+1}u_{n-1} - u_nu_{n-1} + u_{n+1}u_{n-2} + 2(-1)^{n+1},$$
        and hence, using formula (3):
        $$u_n^2 - u_{n+3}u_{n-3} =  - u_nu_{n-1} + u_{n+1}u_{n-2} + 3(-1)^{n+1},$$
        If we now rewrite $u_{n+1} = u_n + u_{n-1}$ and $u_{n-2} = u_n - u_{n-1}$, we get that $u_{n+1}u_{n-2} = u_n^2 - u_{n-1}^2$. Thus,
        $$-u_nu_{n-1} + u_{n+1}u_{n-2} = u_n^2 - u_{n-1}^2 - u_nu_{n-1} = u_n^2 - u_{n+1}u_{n-1} = (-1)^{n+1}$$
        by formula (3), which lets us conclude that
        $$u_n^2 - u_{n+3}u_{n-3} = 4(-1)^{n+1}.$$
        \item Formula (3) tells us that $u_{n+1}u_{n+3} = u_{n+2}^2 + (-1)^n$ and part (c) tells us that $u_nu_{n+4} = u_{n+2}^2 - (-1)^n$. Multiplying both equations gives us 
        $$u_nu_{n+1}u_{n+3}u_{n+4} = (u_{n+2}^2 + (-1)^n)(u_{n+2}^2 - (-1)^n) = u_{n+2}^4 - 1.$$
        \item If we write $u_{2n+2} = u_{2n + 1} + u_{2n}$ and $u_{2n - 1} = u_{2n + 1} - u_{2n}$, then $u_{2n+2}u_{2n-1} = u_{2n+1}^2 - u_{2n}^2$, and hence we obtain
        $$u_{2n+2}u_{2n-1} - u_{2n}u_{2n+1} = u_{2n+1}^2 - u_{2n}^2 - u_{2n}u_{2n+1}.$$
        Next, it is clear that $u_{2n}^2 + u_{2n}u_{2n+1} = u_{2n+2}u_{2n}$ so we get 
        $$u_{2n+2}u_{2n-1} - u_{2n}u_{2n+1} = u_{2n+1}^2 - u_{2n+2}u_{2n}.$$
        Finally, applying formula (3) gives us
        $$u_{2n+2}u_{2n-1} - u_{2n}u_{2n+1} = 1.$$\\
    \end{enumerate}
\end{solution}

\begin{exercise}
    Represent the integers 50, 75, 100, and 125 as sums of distinct Fibonacci numbers. \\
\end{exercise}

\begin{solution}
    The largest Fibonacci number less than 50 is $u_9 = 34$; the largest Fibonacci number smaller than $50 - u_9 = 16$ is $u_7 = 13$; the number $50 - u_9 - u_7$ is equal to 3, a Fibonacci number. Thus, we obtain $50 = 3 + 13 + 34= u_4 + u_7 + u_9$. A similar procedure lets us find that
    \begin{align*}
        75 &= 2 + 5 + 13 + 55 = u_2 + u_5 + u_7 + u_{10}, \\
        100 &= 3 + 8 + 89 = u_4 + u_6 + u_{11}, \\
        125 &= 2 + 34 + 89 = u_2 + u_9 + u_{11}.\\
    \end{align*}
\end{solution}

\begin{exercise}
    Prove that every positive integer can be written as a sum of distinct terms from the sequence $u_2$, $u_3$, $u_4$, ... (that is, the Fibonacci sequence with $u_1$ deleted). \\
\end{exercise}

\begin{solution}
    Let's prove by induction on $n \geq 4$ that every number in the set $\{1, 2, ..., u_n\}$ can be written as a sum of numbers from the set $\{u_2, ..., u_{n-1}\}$, none repeated. When $n = 4$, this follows from the fact that $1 = u_2$, $2 = u_3$ and $3 = u_2 + u_3$.
    
    Next, suppose that the case $n = k$ holds for some integer $k \geq 4$, i.e., every elements between 1 and $u_k$ can be written as a sum of distinct numbers from the set $\{u_2, u_3, ..., u_{k-1}\}$. Consider an arbitrary integer $u_{k} + 1 \leq N \leq u_{k+1}$, then $1 \leq N- u_k \leq u_{k+1} - u_k = u_{k-1} \leq u_k$. By our hypothesis, $N-u_k$ can be written as a sum of distinct elements from the set $\{u_2, ..., u_{k-1}\}$, so it directly follows that $N$ can be written as a sum of distinct elements from the set $\{u_2, ..., u_k\}$. This proves that the case $n = k+1$ holds as well. Thus, it holds for all $n \geq 4$ by induction.
    
    Therefore, every  positive integer can be written as a sum of distinct terms from the sequence $u_2$, $u_3$, $u_4$, ... since $u_n \rightarrow \infty$ as $n \rightarrow \infty$. \\
\end{solution}

\begin{exercise}
    Establish the identity
    $$(u_nu_{n+3})^2 + (2u_{n+1}u_{n+2})^2 = (u_{2n+3})^2, \qquad \qquad \qquad n \geq 1$$
    and use this to generate five primitive Pythagorean triples. \\
\end{exercise}

\begin{solution}
    If we write $u_n = u_{n + 2} - u_{n-1}$ and $u_{n+3} = u_{n+2} + u_{n+1}$, then $u_nu_{n+3} = u_{n+2}^2 - u_{n+1}^2$. Thus,
    $$(u_nu_{n+3})^2 = u_{n+2}^4 - 2u_{n+1}^2u_{n-1}^2 + u_{n+1}^4,$$
    and so
    $$(u_nu_{n+3})^2 + (2u_{n+1}u_{n-1})^2 = u_{n+2}^4 + 2u_{n+1}^2u_{n-1}^2 + u_{n+1}^4 = (u_{n+2}^2 + u_{n+1}^2)^2.$$
    Finally, by Problem 3, we obtain
    $$(u_nu_{n+3})^2 + (2u_{n+1}u_{n-1})^2 = (u_{2n + 3})^2.$$
    Let's now generate five primitive Pythagorean triples. If we take $n = 1$, we obtain the following primitive Pythagorean triple:
    $$3^2 + 4^2 = 5^2.$$
    If we take $n = 2$, we obtain the following primitive Pythagorean triple:
    $$5^2 + 12^2 = 13^2.$$
    If we take $n = 3$, we obtain the triple $16^2 + 30^2 = 34^2$ which induces the following primitive Pythagorean triple:
    $$8^2 + 15^2 = 15^2.$$
    If we take $n = 4$, we obtain the following primitive Pythagorean triple:
    $$39^2 + 80^2 = 89^2.$$
    Finally, if we take $n = 5$, we obtain the following primitive Pythagorean triple:
    $$105^2 + 208^2 = 233^2.$$ \\
\end{solution}

\begin{exercise}
    Prove that the product $u_nu_{n+1}u_{n+2}u_{n+3}$ of any four consecutive Fibonacci numbers is the area of a Pythagorean triangle. \hint{See the previous problem} \\
\end{exercise}

\begin{solution}
    By the previous problem, we have that $u_nu_{n+3}$, $2u_{n+1}u_{n+2}$, $u_{2n+3}$ are the sides of a right triangle. Thus, it directly follows that its area is 
    $$\frac{(u_nu_{n+3})(2u_{n+1}u_{n+2})}{2} = u_nu_{n+1}u_{n+2}u_{n+3}.$$ \\
\end{solution}

\begin{exercise}
    Let $\alpha = \frac{1}{2}(1+ \sqrt{5})$ and $\beta = \frac{1}{2}(1 - \sqrt{5})$, so that $\alpha$ and $\beta$ are both roots of the equation $x^2 = x + 1$. Show by induction that the \textit{Binet formula}
    $$u_n = \frac{\alpha^n - \beta^n}{\sqrt{5}}$$
    holds for $n \geq 1$. \\
\end{exercise}

\begin{solution}
    The case $n = 1$ follows from the fact that 
    $$\frac{\alpha^1 - \beta^1}{\sqrt{5}} = \frac{\sqrt{5}}{\sqrt{5}} = 1 = u_1.$$
    For the case $n = 2$, recall that both $\alpha$ and $\beta$ are solutions of the equation $x^2 = x + 1$, i.e., $\alpha^2 = \alpha + 1$ and $\beta^2 = \beta + 1$. From this, we have that $\alpha^2 - \beta^2 = \sqrt{5}$. Thus, we get 
    $$\frac{\alpha^2 - \beta^2}{\sqrt{5}} = \frac{\sqrt{5}}{\sqrt{5}} = 1 = u_2.$$
    Next, suppose that the cases $n = k-1$ and $n = k$ hold for some $k \geq 2$. Since $\alpha^2 = \alpha + 1$ and $\beta^2 = \beta + 1$, then $\alpha^{k+1} = \alpha^k + \alpha^{k-1}$ and $\beta^{k+1} = \beta^k + \beta^{k-1}$. Therefore:
    \begin{align*}
        u_{k+1} &= u_k + u_{k-1} \\
        &= \frac{\alpha^k - \beta^k}{\sqrt{5}} + \frac{\alpha^{k-1} - \beta^{k-1}}{\sqrt{5}} \\
        &= \frac{(\alpha^k + \alpha^{k-1}) - (\beta^k + \beta^{k-1})}{\sqrt{5}} \\
        &= \frac{\alpha^{k+1} - \beta^{k+1}}{\sqrt{5}}.
    \end{align*}
    Hence, the case $n = k+1$ holds as well, and so it holds for all $n \geq 1$ by induction. \\
\end{solution}

\begin{exercise}
    For $n \geq 1$, show that the product $u_{2n-1}u_{2n+5}$ can be expressed as the sum of two squares. \hint{Problem 6(e)} \\
\end{exercise}

\begin{solution}
    When $n = 1$, we have $u_{2n-1}u_{2n+5} = 1\cdot 13 = 2^2 + 3^2$. When $n \geq 2$, Problem 3 tells us that both $u_{2n-1}$ and $u_{2n+5}$ can be expressed as the sum of two squares, and hence, their product too. More explicitly,
    $$u_{2n-1}u_{2n+5} = (u_n^2 + u_{n-1}^2)(u_{n+3}^2 + u_{n+2}^2) = (u_nu_{n+3} + u_{n-1}u_{n+2})^2 + (u_nu_{n+2} - u_{n-1}u_{n+3})^2.$$ \\
\end{solution}

\begin{exercise}
    Prove that if $p = 4k + 3$ is prime, then $p$ cannot divide a Fibonacci number with odd index; that is, $p \not\mid u_{2n-1}$ for all $n \geq 1$. \hint{In the contrary case, $u_n^2 + u_{n-1}^2 = u_{2n-1} \equiv 0 \Mod{p}$. See Problem 12, Section 5.4} \\
\end{exercise}

\begin{solution}
    Suppose that $p \mid u_{2n-1}$ for some $n \geq 1$. When $n = 1$, we have $u_{2n-1} = 1$ so this is impossible. When $n \geq 2$, we have that $u_{2n-1} = u_n^2 + u_{n-1}^2$. Thus, $u_n^2 + u_{n-1}^2 \equiv 0 \Mod{p}$. If $p \mid u_n$, then the congruence becomes $u_{n-1}^2 \equiv 0 \Mod{p}$ which implies that $p \mid u_{n-1}$. But this is impossible since $\gcd(u_n, u_{n-1}) = 1$. Since $\gcd(u_n, p) = 1$, there is an integer $a$ such that the congruence $au_n \equiv 1 \Mod{p}$. Hence, the congruence $u_n^2 + u_{n-1}^2 \equiv 0 \Mod{p}$ implies that $(au_{n-1})^2 \equiv -1 \Mod{p}$, and hence, $p \equiv 1 \Mod{4}$. Since this clearly contradicts the fact that $p = 4k + 3$, then $p$ can't divide $u_{2n-1}$. \\ 
\end{solution}

\begin{exercise}
    Verify that the product $u_{2n}u_{2n+2}u_{2n+4}$ of three consecutive Fibonacci numbers with even indices is the product of three consecutive integers; for instance, $u_4u_5u_6 = 504 = 7\cdot 8 \cdot 9$. \hint{First show that $u_{2n}u_{2n+4} = u_{2n+2}^2 - 1$} \\
\end{exercise}

\begin{solution}
    By Problem 6(c), we have $u_{2n}u_{2n+4} = u_{2n+2}^2 - 1 = (u_{2n+2} - 1)(u_{2n+2} + 1)$. Therefore,
    $$u_{2n}u_{2n+2}u_{2n+4} = (u_{2n+2} - 1)u_{2n+2}(u_{2n+2} + 1).$$ \\
\end{solution}

\begin{exercise}
    Use formulas (1) and (2) to show that the sum of any twenty consecutive Fibonacci numbers is divisible by $u_{10}$. \\
\end{exercise}

\begin{solution}
    By formula (1), we have that
    $$\sum_{k=0}^{19}u_{n+k} = \sum_{k=1}^{n+19}u_k - \sum_{k=1}^{n-1}u_k = u_{n + 21} - u_{n+1}.$$
    Next, formula (2) tells us that $u_{n+21} = u_{20}u_n + u_{21}u_{n+1}$. So it follows that
    $$u_n + u_{n+1} + ... + u_{n+19} = u_{20}u_n + (u_{21} - 1)u_{n+1}.$$
    It is clear that $u_{20}$ is divisible by $u_{10}$ so it suffices to show that $u_{10} = 5$ divides $u_{21} - 1 = 10945$ as well. Since it is clearly divisible by 5 and $10945 = 11\cdot 995$, then $u_{10} \mid u_{21} - 1$. Therefore, $u_{10}$ divides $u_{20}u_n + (u_{21} - 1)u_{n+1} = \sum_{k=0}^{19}u_{n+19}$.  \\
\end{solution}

\begin{exercise}
    For $n \geq 4$, prove that $u_n + 1$ is not a prime. \hint{It suffices to establish the identities
    \begin{align*}
        u_{4k} + 1 &= u_{2k - 1}(u_{2k} + u_{2k+2}) \\
        u_{4k+1} + 1 &= u_{2k + 1}(u_{2k-1} + u_{2k+1}) \\
        u_{4k+2} + 1 &= u_{2k + 2}(u_{2k+1} + u_{2k-1}) \\
        u_{4k+3} + 1 &= u_{2k + 1}(u_{2k+1} + u_{2k+3})} \\
    \end{align*}
\end{exercise}

\begin{solution}
    Let $k \geq 1$, then formula (1) tells us that 
    $$u_{4k} + 1 = u_{2k - 1}u_{2k} + u_{2k + 1}u_{2k} + 1.$$
    Since $u_{2k + 1} = u_{2k + 2} - u_{2k}$, then we can rewrite the above equation as follows:
    $$u_{4k} + 1 = u_{2k - 1}u_{2k} + u_{2k + 2}u_{2k} - u_{2k}^2 + 1.$$
    Now, if we simply recall that $u_{2k + 2}u_{2k} + 1 = u_{2k + 1}^2$ by formula (3), then
    $$u_{4k} + 1 = u_{2k - 1}u_{2k} + u_{2k + 1}^2 - u_{2k}^2.$$
    Finally, $u_{2k + 1}^2 - u_{2k}^2 = (u_{2k + 1} - u_{2k})(u_{2k + 1} + u_{2k}) = u_{2k - 1}u_{2k+2}$ so
    $$u_{4k} + 1 = u_{2k - 1}(u_{2k} + u_{2k+2}).$$
    Therefore, if $n \equiv 0 \Mod{4}$, $u_n + 1$ can't be prime since $u_{(n/2) - 1} \mid u_n + 1$ and $u_{(n/2) - 1} > 1$.\\
    
    Similarly, formula (1) gives us
    $$u_{4k + 1} + 1 = u_{2k + 1}u_{2k - 1} + u_{2k + 2}u_{2k} + 1.$$
    By formula (3), $u_{2k + 2}u_{2k} + 1 = u_{2k + 1}^2$:
    $$u_{4k + 1} + 1 = u_{2k + 1}u_{2k - 1} + u_{2k+1}^2 = u_{2k + 1}(u_{2k - 1} + u_{2k+1}).$$
    Therefore, if $n \equiv 1 \Mod{4}$, $u_n + 1$ can't be prime since $u_{(n+1)/2} \mid u_n + 1$ and $u_{(n+1)/2} > 1$.\\
    
    Next, formula (1) tells us that 
    $$u_{4k + 2} + 1 = u_{2k + 1}u_{2k} + u_{2k + 2}u_{2k + 1} + 1.$$
    Since $u_{2k + 1} = u_{2k + 2} - u_{2k}$, then we can rewrite the above equation as follows:
    $$u_{4k + 2} + 1 = u_{2k + 2}u_{2k + 1} + u_{2k + 2}u_{2k} - u_{2k}^2 + 1.$$
    By formula (3), $u_{2k + 2}u_{2k} + 1 = u_{2k + 1}^2$ so we get
    $$u_{4k + 2} + 1 = u_{2k + 2}u_{2k + 1} + u_{2k + 1}^2 - u_{2k}^2.$$
    Finally, $u_{2k + 1}^2 - u_{2k}^2 = (u_{2k + 1} - u_{2k})(u_{2k + 1} + u_{2k}) = u_{2k - 1}u_{2k+2}$ so
    $$u_{4k + 2} + 1 = u_{2k + 2}(u_{2k+1} + u_{2k-1}).$$
    Therefore, if $n \equiv 2 \Mod{4}$, $u_n + 1$ can't be prime since $u_{(n + 2)/2} \mid u_n + 1$ and $u_{(n + 2)/2} > 1$. \\

    Lastly, formula (1) gives us
    $$u_{4k + 3} + 1 = u_{2k}u_{2k + 2} + u_{2k + 1}u_{2k + 3} + 1.$$
    By formula (3), $u_{2k + 2}u_{2k} + 1 = u_{2k + 1}^2$:
    $$u_{4k + 3} + 1 = u_{2k + 1}(u_{2k + 1} + u_{2k+3}).$$
    Therefore, if $n \equiv 3 \Mod{4}$, $u_n + 1$ can't be prime since $u_{(n-1)/2} \mid u_n + 1$ and $u_{(n-1)/2} > 1$. \\
\end{solution}

\begin{exercise}
    The \textit{Lucas numbers} are defined by the same recurrence formula as the Fibonacci numbers,
    $$L_n = L_{n-1} + L_{n-2} \quad (n \geq 3),$$
    but with $L_1 = 1$ and $L_2 = 3$; this gives the sequence 1, 3, 4, 7, 11, 18, 29, 47, 76, 123, 199, 322, ... For the Lucas numbers, derive each of the identities below:
    \begin{enumerate}
        \item $L_1 + L_2 + L_3 + ... + L_n = L_{n+2} - 3$, $n \geq 1$.
        \item $L_1 + L_3 + L_5 + ... + L_{2n-1} = L_{2n} - 2$, $n \geq 1$.
        \item $L_2 + L_4 + L_4 + ...  + L_{2n} = L_{2n+1} - 1$, $n \geq 1$.
        \item $L_n^2 = L_{n+1}L_{n-1} + 5(-1)^n$, $n \geq 2$.
        \item $L_1^2 + L_2^2 + L_3^2 + ... + L_n^2 = L_nL_{n+1} - 2$, $n \geq 1$.
        \item $L_{n+1}^2 - L_n^2 = L_{n-1}L_{n+2}$, $n \geq 1$.\\
    \end{enumerate}
\end{exercise}

\begin{solution}
    \begin{enumerate}
        \item Since $L_1 = L_3 - L_2$, $L_2 = L_4 - L_3$, $L_3 = L_5 - L_4$, ..., $L_n = L_{n+2} - L_{n+1}$, then adding all these equations together gives us $\sum_{k=1}^{n}L_k = L_{n+2} - L_2$ since the right side is telescopic. Since $L_2 = 3$, then we get 
        $$L_1 + L_2 + L_3 + ... + L_n = L_{n+2} - 3.$$
        \item Since $L_1 = 1$, $L_3 = L_4 - L_2$, $L_5 = L_6 - L_4$, ..., $L_{2n-1} = L_{2n} - L_{2n-2}$, then adding all these equations together gives us $\sum_{k=1}^{n}L_{2k-1} = L_{2n} - L_2 + 1$ since the right side is telescopic. Since $L_2 = 3$, then we get 
        $$L_1 + L_3 + L_5 + ... + L_{2n-1} = L_{2n} - 2.$$
        \item If we let $S = \sum_{k=1}^{n}L_{2k}$, $S_1 = \sum_{k=1}^{n}L_{2k-1}$ and $S_2 = \sum_{k=1}^{2n}L_k$, then $S = S_2 - S_1$. Moreover, part (a) and part (b) imply that $S_1 = L_{2n} - 2$ and $S_2 = L_{2n+2} - 3$, so we get that $S = (L_{2n+2} - 3) - (L_{2n} - 2)$. Since $(L_{2n+2} - 3) - (L_{2n} - 2) = L_{2n+1} - 1$, then 
        $$L_2 + L_4 + L_4 + ...  + L_{2n} = L_{2n+1} - 1.$$
        \item Let's prove it by induction on $n \geq 2$. For the case $n = 2$, we have that both $L_2^2$ and $L_3L_1 + 5(-1)^2$ are equal to 9, and so it holds for that case. Next, suppose that there is an integer $k \geq 2$ such that it holds for $n = k$, then $L_k^2 - L_{k+1}L_{k-1} = 5(-1)^k$. It follows that
        $$L_{k+1}^2 - L_{k+2}L_k = L_{k+1}^2 - L_{k+1}L_k - L_k^2 = -L_{k-1}L_{k+1} - L_k^2 = 5(-1)^{k+1},$$
        and so it holds for $n = k+1$ as well. Therefore, it holds for all $n \geq 2$ by induction.
        \item Let's prove it by induction on $n \geq 1$. The case $n = 1$ holds because both $L_1^2$ and $L_1L_2 - 2$ are equal to 1. Next, suppose that it holds for some integer $k \geq 1$, then
        $$L_1^2 + L_2^2 + ... + L_k^2 = L_kL_{k+1} - 2.$$
        Adding $L_{k+1}^2$ on both sides gives us 
        $$L_1^2 + L_2^2 + ... + L_{k+1}^2 = L_kL_{k+1} + L_{k+1}^2 - 2.$$
        Since $L_kL_{k+1} + L_{k+1}^2 = (L_k + L_{k+1})L_{k+1} = L_{k+1}L_{k+2}$, then 
        $$L_1^2 + L_2^2 + ... + L_{k+1}^2 = L_{k+1}L_{k+2} - 2$$
        which is exactly the case $n = k+1$. Therefore, it holds for all $n \geq 1$ by induction.
        \item By direct calculation:
        $$L_{n+1}^2 - L_n^2 = (L_{n+1} - L_n)(L_{n+1} + L_n) = L_{n-1}L_{n+2}.$$ \\
    \end{enumerate}
\end{solution}

\begin{exercise}
    Establish the following relations between the Fibonacci and Lucas numbers:
    \begin{enumerate}
        \item $L_n = u_{n+1} + u_{n-1} = u_{n} + 2u_{n-1}$, $n \geq 2$. \hint{Argue by induction on $n$}
        \item $L_n = u_{n+2} - u_{n-2}$, $n \geq 3$.
        \item $u_{2n} = u_nL_n$, $n \geq 1$.
        \item $L_{n+1} + L_{n-1} = 5u_n$, $n \geq 1$.
        \item $L_n^2 = u_n^2 + 4u_{n+1}u_{n-1}$, $n \geq 2$.
        \item $2u_{m+n} = u_mL_n + L_mu_n$, $m \geq 1$, $n\geq1$.
        \item $\gcd(u_n, L_n) = 1$ or 2, $n \geq 1$.\\
    \end{enumerate}
\end{exercise}

\begin{solution}
    \begin{enumerate}
        \item Let's prove it by induction. When $n = 2$, both $L_2$ and $u_3 + u_1$ are equal to 3 so it holds. When $n = 3$, both $L_3$ and $u_4 + u_2$ are equal to 4 so it holds as well. Next, suppose there is an integer $k \geq 3$ such that it holds for $n = k-1$ and $n = k$, then we get the equations $L_{k-1} = u_k + u_{k-2}$ and $L_k = u_{k+1} + u_{k-1}$. Adding these equations together gives
        $$L_{k+1} = L_k + L_{k-1} = (u_{k+1} + u_k) + (u_{k-1} + u_{k-2}) = u_{k+2} + u_k,$$
        implying that the case $n = k+1$ holds. Therefore, it holds for all $n \geq 2$ by induction. Finally, $u_{n+1} = u_n + u_{n-1}$ so $L_n = u_{n+1} + u_{n-1} = u_n + 2u_{n-1}$.
        \item Since $u_{n+2} = u_{n+1} + u_n$ and $u_n = u_{n-1} + u_{n-2}$, then $u_{n+1} + u_{n-1} =u_{n+2} - u_{n-2}$. Therefore, using part (a), $L_n = u_{n+2} - u_{n-2}$ for all $n \geq 3$.
        \item By formula (1), $u_{2n} = u_{n-1}u_n + u_nu_{n+1} = u_n(u_{n-1} + u_{n+1})$. Hence, by part (a), $u_{2n} = u_nL_n$ for all $n \geq 1$.
        \item By part (a), $L_{n+1} + L_{n-1} = u_{n+2} + u_n + u_n + u_{n-2}$. Since $u_{n+2} = u_{n+1} + u_n$ and $u_{n+1} = u_n + u_{n-1}$, then $u_{n+2} + u_{n-2} = 3u_n$. It follows that $L_{n+1} + L_{n-1} = 5u_n$ for all $\geq 1$.
        \item If square the formula in part (a), we get $L_n^2 = u_{n+1}^2 + u_{n-1}^2 + 2u_{n+1}u_{n-1}$. Now, noitice that squaring both sides of the identity $u_n = u_{n+1} - u_{n-1}$ gives us that $u_{n+1}^2 + u_{n-1}^2 = u_n^2 + 2u_{n+1}u_{n-1}$. Therefore, $L_n^2 = u_n^2 + 4u_{n+1}u_{n-1}$ for all $n \geq 2$.
        \item Formula (1) gives us that $u_{m+n} = u_{m-1}u_n + u_mu_{n+1}$ and $u_{m+n} = u_{n-1}u_m + u_nu_{m+1}$ if we interchange $m$ and $n$. Adding these two equations gives us $2u_{m+n} = u_m(u_{n-1} + u_{n+2}) + (u_{m-1} + u_{m+1})u_n$. Therefore, by part (a), $2u_{m+n} = u_mL_n + L_mu_n$ for all $m,n \geq 2$. 
        
        Notice that part (a) is only valid when $m,n \geq 2$ so the cases $m = 1$ and $n = 1$ have to be proved independently. Since the formula is symmetric, it suffices to prove the case $m = 1$, $n \geq 2$. Clearly:
        $$u_1L_n + L_1u_n = u_{n+1} + u_{n-1} + u_n = 2u_{n+1}.$$
        Therefore, the formula holds for all $m,n \geq 1$.
        \item When $n = 1$, this is trivial since both are equal to 1. When $n \geq 2$, let $d$ be a common positive divisor of $u_n$ and $L_n$. Since part (a) tells us that $L_n = u_n + 2u_{n-1}$, then $d \mid L_n - u_n = 2u_{n-1}$. Since $\gcd(u_n, u_{n-1}) = 1$, then $\gcd(d, u_{n-1}) = 1$, so $d \mid 2u_{n-1}$ implies that $d \mid 2$. Therefore, $\gcd(u_n, L_n) = 1$ or 2 for all $n \geq 1$. \\
    \end{enumerate}
\end{solution}

\begin{exercise}
    If $\alpha = \frac{1}{2}(1 + \sqrt{5})$ and $\beta = \frac{1}{2}(1 - \sqrt{5})$, prove that 
    $$L_n = \alpha^n + \beta^n$$
    for $n \geq 1$; whence
    $$\alpha^n = \frac{1}{2}(L_n + \sqrt{5}u_n), \quad \beta^n = \frac{1}{2}(L_n - \sqrt{5}u_n).$$
\end{exercise}

\begin{solution}
    Let's prove it by induction on $n$. The case $n = 1$ follows from the fact that 
    $$\alpha^1 + \beta^1 = 1 = L_1.$$
    For the case $n = 2$, recall that both $\alpha$ and $\beta$ are solutions of the equation $x^2 = x + 1$, i.e., $\alpha^2 = \alpha + 1$ and $\beta^2 = \beta + 1$. From this, we have that 
    $$\alpha^2 + \beta^2 = \alpha + \beta + 1 + 1 = 3 = L_2.$$
    Next, suppose that the cases $n = k-1$ and $n = k$ hold for some $k \geq 2$. Since $\alpha^2 = \alpha + 1$ and $\beta^2 = \beta + 1$, then $\alpha^{k+1} = \alpha^k + \alpha^{k-1}$ and $\beta^{k+1} = \beta^k + \beta^{k-1}$. Therefore:
    $$L_{k+1} = L_k + L_{k-1} = (\alpha^k + \alpha^{k-1}) + (\beta^k + \beta^{k-1}) = \alpha^{k+1} + \beta^{k+1}.$$
    Hence, the case $n = k+1$ holds as well, and so it holds for all $n \geq 1$ by induction.\\
\end{solution}

\begin{exercise}
    In 1876, Lucas discovered the following formula for the Fibonacci numbers in terms of the binomial coefficients:
    $$u_n = \binom{n-1}{0} + \binom{n-2}{1} +\binom{n-3}{2} + \dots + \binom{n-j}{j-1} + \binom{n-j-1}{j}$$
    where $j$ is the largest integer less than or equal to $(n-1)/2$. Derive this result. \hint{Argue by induction, using the relation $u_n = u_{n-1} + u_{n-2}$; note also that $\binom{m}{k} = \binom{m-1}{k} + \binom{m-1}{k-1}$} \\
\end{exercise}

\begin{solution}
    Let's prove it by induction on $n$. When $n = 1$, we have that $j = 0$ so $\sum_{k=0}^{j}\binom{n-k-1}{k} = \binom{0}{0} = 1 = u_1$. When $n = 2$, we have that $j = 0$ as well so $\sum_{k=0}^{j}\binom{2-k-1}{k} = \binom{1}{0} = 1 = u_2$. Now, suppose that the formula holds for $n = k-1$ and $n = k$ for some $k \geq 2$, and denote by $j_i$ the largest integer less than or equal to $(i-1)/2$ for $i = k-1, k, k+1$. Let's first consider the case where $k$ is even, in that case, $j_{k_1} = j_k = (k/2) - 1$ and $j_{k+1} = k/2$. Hence, 
    $$u_{k-1} = \sum_{i=0}^{(k/2) - 1}\binom{k-i-2}{i},\quad \text{ and } \quad u_k = \sum_{i=0}^{(k/2) - 1}\binom{k-i-1}{i}.$$
    If we now reindex the sums in the following way
    $$u_{k-1} = \sum_{i=1}^{(k/2) - 1}\binom{k-i-1}{i-1} + 1,\qquad u_k = 1 + \sum_{i=1}^{(k/2) - 1}\binom{k-i-1}{i}$$
    and add these two equations together, we obtain 
    $$u_{k+1} = 1 + \sum_{i=1}^{(k/2) - 1}\left[\binom{k-i-1}{i-1} + \binom{k-i-1}{i}\right] + 1.$$
    The summation formula for the binomial coefficients tells us that $\binom{k-i-1}{i-1} + \binom{k-i-1}{i} = \binom{k-i}{i}$. Moreover, $\binom{k}{0} = \binom{k - j_{k+1}}{j_{k+1}} = 1$ which gives us
    $$u_{k+1} = \binom{k}{0} + \sum_{i=1}^{(k/2) - 1}\binom{k-i}{i-1} + \binom{k-j_{k+1}}{j_{k+1}} = \sum_{i=0}^{j_{k+1}}\binom{(k+1) - i - 1}{i}.$$
    Thus, the case $n = k+1$ holds. 
    
    Next, let's now suppose that $k$ is odd, then $j_{k-1} = (k-1)/2 - 1$ and $j_k = j_{k+1} = (k-1)/2$. Hence:
    $$u_{k-1} = \sum_{i=0}^{(k-1)/2 - 1}\binom{k-i-2}{i}, \quad \text{ and } \quad u_k = \sum_{i=0}^{(k-1)/2}\binom{k-i-1}{i}.$$
    If we now reindex the sums in the following way
    $$u_{k-1} = \sum_{i=1}^{(k-1)/2}\binom{k-i-1}{i-1}, \quad \text{ and } \quad u_k = 1 + \sum_{i=1}^{(k-1)/2}\binom{k-i-1}{i}.$$
    and add these two equations together, we obtain 
    $$u_{k+1} = 1 + \sum_{i=1}^{(k-1)/2}\left[\binom{k-i-1}{i-1} + \binom{k-i-1}{i}\right].$$
    The summation formula for the binomial coefficients tells us that $\binom{k-i-1}{i-1} + \binom{k-i-1}{i} = \binom{k-i}{i}$. Moreover, $\binom{k}{0} = 1$ which gives us
    $$u_{k+1} = \binom{k}{0} + \sum_{i=1}^{(k-1)/2}\binom{k-i}{i-1} = \sum_{i=0}^{j_{k+1}}\binom{(k+1) - i - 1}{i}.$$
    Thus, the case $n = k+1$ holds again. Therefore, it holds for all $n \geq 1$ by induction. \\
\end{solution}

\begin{exercise}
    Establish that for $n \geq 1$,
    \begin{enumerate}
        \item $\binom{n}{1}u_1 + \binom{n}{2}u_2 + \binom{n}{3}u_3 + \dots + \binom{n}{n}u_n = u_{2n}$;
        \item $-\binom{n}{1}u_1 + \binom{n}{2}u_2 - \binom{n}{3}u_3 + \dots + (-1)^n\binom{n}{n}u_n = -u_n$.\\
    \end{enumerate}
\end{exercise}

\begin{solution}
    \begin{enumerate}
        \item By the Binet formula:
        $$\sum_{k=1}^{n}\binom{n}{k}u_k = \frac{1}{\sqrt{5}}\left[\sum_{k=1}^{n}\binom{n}{k}\alpha^k - \sum_{k=1}^{n}\binom{n}{k}\beta^k\right].$$
        The term $k = 0$ in both sums is 1 so it would change nothing to make both sums begin at that term since it would cancel out anyway:
        $$\sum_{k=1}^{n}\binom{n}{k}u_k = \frac{1}{\sqrt{5}}\left[\sum_{k=0}^{n}\binom{n}{k}\alpha^k - \sum_{k=0}^{n}\binom{n}{k}\beta^k\right].$$
        Next, by the Binomial Theorem:
        $$\sum_{k=1}^{n}\binom{n}{k}u_k = \frac{1}{\sqrt{5}}\left[(\alpha + 1)^n - (\beta + 1)^n\right].$$
        Recall now that $\alpha + 1 = \alpha^2$ and $\beta + 1 = \beta^2$:
        $$\sum_{k=1}^{n}\binom{n}{k}u_k = \frac{\alpha^{2n} - \beta^{2n}}{\sqrt{5}}.$$
        Finally, using the Binet formula again gives us 
        $$\binom{n}{1}u_1 + \binom{n}{2}u_2 + \binom{n}{3}u_3 + \dots + \binom{n}{n}u_n = u_{2n}.$$
        \item By the Binet formula:
        $$\sum_{k=1}^{n}\binom{n}{k}(-1)^ku_k = \frac{1}{\sqrt{5}}\left[\sum_{k=1}^{n}\binom{n}{k}(-1)^k\alpha^k - \sum_{k=1}^{n}\binom{n}{k}(-1)^k\beta^k\right].$$
        The term $k = 0$ in both sums is 1 so it would change nothing to make both sums begin at that term since it would cancel out anyway:
        $$\sum_{k=1}^{n}\binom{n}{k}(-1)^ku_k = \frac{1}{\sqrt{5}}\left[\sum_{k=0}^{n}\binom{n}{k}(-1)^k\alpha^k - \sum_{k=0}^{n}\binom{n}{k}(-1)^k\beta^k\right].$$
        Next, by the Binomial Theorem:
        $$\sum_{k=1}^{n}\binom{n}{k}(-1)^ku_k = \frac{1}{\sqrt{5}}\left[(1 - \alpha)^n - (1 - \beta)^n\right].$$
        Recall now that $\alpha + 1 = \alpha^2$ and $\beta + 1 = \beta^2$, and so $1 - \alpha = -1/\alpha$ and $1 - \beta = 1/\beta$:
        $$\sum_{k=1}^{n}\binom{n}{k}(-1)^ku_k = \frac{\alpha^{-n} - \beta^{-n}}{\sqrt{5}} = -\frac{(-1)^n}{\alpha^n\beta^n}\cdot\frac{\alpha^n - \beta^n}{\sqrt{5}}.$$
        Finally, notice that $\alpha\beta = -1$ and use the Binet formula again to obtain:
        $$-\binom{n}{1}u_1 + \binom{n}{2}u_2 - \binom{n}{3}u_3 + \dots + (-1)^n\binom{n}{n}u_n = -u_n.$$ \\
    \end{enumerate}
\end{solution}

\begin{exercise}
    Prove that 24 divides the sum of any 24 consecutive Fibonacci numbers. \hint{Consider the identity
    $$u_n + u_{n+1} + ... + u_{n + k -1} = u_{n-1}(u_{k+1} - 1) + u_n(u_{k+2} - 1)}$$
\end{exercise}

\begin{solution}
    First, notice that
    $$u_n + u_{n+1} + ... + u_{n+23} = \sum_{k=1}^{n+23}u_k - \sum_{k=1}^{n-1}u_k.$$
    By formula (2), $\sum_{k=1}^{n+23}u_k = u_{n+25} - 1$ and $\sum_{k=1}^{n-1}u_k = u_{n+1} -1$ so
    $$u_n + u_{n+1} + ... + u_{n+23} = u_{n+25} - u_{n+1}.$$
    Using formula (1) to get $u_{n+25} = u_{n-1}u_{25} + u_nu_{26}$ and $u_{n+1} = u_n + u_{n-1}$ gives us
    $$u_n + u_{n+1} + ... + u_{n+23} = u_{n-1}(u_{25} - 1) + u_n(u_{26} - 1).$$
    Finally, it suffices to see that $u_{25} - 1 = 24\cdot 3126$ and $u_{26} - 1 = 24\cdot 5058$ to conclude that 24 divides the sum on the left hand side. Therefore, 24 divides any sum of 24 consecutive Fibonacci numbers. \\
\end{solution}

\begin{exercise}
    Let $n$ be a positive integer and $m = n^{13} - n$. Show that $u_m$ is divisible by 30290. \hint{See Problem 1(b) of Section 7.3} \\
\end{exercise}

\begin{solution}
    By Fermat's Little Theorem, $a^2 \equiv 1 \Mod{3}$ for all integers $a$ coprime with 3. Taking both sides to the sixth power and multiplying both sides by $a$ gives us that $a^{13}\equiv a \Mod{3}$ holds for all integers $a$. Thus, in the particular case $a = n$, we get that $3 \mid n^{13} - n = m$. It follows that $u_3 = 2 \mid u_m$.

    By Fermat's Little Theorem, $a^4 \equiv 1 \Mod{5}$ for all integers $a$ coprime with 5. Taking both sides to the third power and multiplying both sides by $a$ gives us that $a^{13}\equiv a \Mod{5}$ holds for all integers $a$. Thus, in the particular case $a = n$, we get that $5 \mid n^{13} - n = m$. It follows that $u_5 = 5 \mid u_m$.

    By Fermat's Little Theorem, $a^6 \equiv 1 \Mod{7}$ for all integers $a$ coprime with 7. Squaring and multiplying both sides by $a$ gives us that $a^{13}\equiv a \Mod{7}$ holds for all integers $a$. Thus, in the particular case $a = n$, we get that $7 \mid n^{13} - n = m$. It follows that $u_7 = 13 \mid u_m$.

    By Fermat's Little Theorem, $a^{12} \equiv 1 \Mod{13}$ for all integers $a$ coprime with 13. Multiplying both sides by $a$ gives us that $a^{13}\equiv a \Mod{13}$ holds for all integers $a$. Thus, in the particular case $a = n$, we get that $13 \mid n^{13} - n = m$. It follows that $u_{13} = 233 \mid u_m$.

    Finally, 2, 5, 13 and 233 are pairwise relatively prime and they all divide $u_m$. Therefore, it follows that $30290 = 2\cdot 5 \cdot 13 \cdot 233 \mid u_m$.
\end{solution}